%============================================================
%  LEX v3.0 — Template Visualisasi Akademik
%  Perintah Copilot: /generate-viz [jenis data]
%  Tersedia: Bar Chart, Pie Chart, Radar Chart, Tabel Perbandingan
%============================================================

\documentclass[12pt,a4paper]{article}
\usepackage[margin=2.5cm]{geometry}
\usepackage{tikz}
\usetikzlibrary{positioning}
\usepackage{pgfplots}
\pgfplotsset{compat=1.18}
\usepackage{pgf-pie}
\usepackage{booktabs}
\usepackage{multirow}
\usepackage{array}
\usepackage[bahasa]{babel}

\begin{document}

%------------------------------------------------------------
%  1. BAR CHART — Statistik Putusan / Data Empiris
%     Perintah: /generate-viz statistik-putusan
%     Ganti data di \addplot coordinates dengan data riil Anda
%------------------------------------------------------------
\begin{figure}[h]
\centering
\begin{tikzpicture}
\begin{axis}[
  ybar,
  bar width        = 0.5cm,
  width            = 11cm,
  height           = 6cm,
  enlarge x limits = 0.15,
  ylabel           = {Jumlah Putusan},
  xlabel           = {Tahun},
  symbolic x coords= {2019, 2020, 2021, 2022, 2023},
  xtick            = data,
  nodes near coords,
  nodes near coords align = {vertical},
  ymin=0,
  legend style     = {at={(0.5,-0.18)}, anchor=north, legend columns=-1},
  grid             = major,
  grid style       = dashed,
]
  % Ganti nilai-nilai di bawah dengan data riil Anda
  \addplot+[fill=blue!40]   coordinates {(2019,12)(2020,18)(2021,24)(2022,31)(2023,27)};
  \addplot+[fill=orange!50] coordinates {(2019,8) (2020,11)(2021,15)(2022,19)(2023,22)};
  \legend{[Kategori A], [Kategori B]}
\end{axis}
\end{tikzpicture}
\caption{[Judul Grafik --- sumber: [PERLU VERIFIKASI]]}
\label{fig:statistik-putusan}
\end{figure}

%------------------------------------------------------------
%  2. PIE CHART — Komposisi Data
%     Perintah: /generate-viz komposisi-data
%------------------------------------------------------------
\begin{figure}[h]
\centering
\begin{tikzpicture}
  % Ganti persentase dan label sesuai data Anda (total harus 100)
  \pie[
    text         = legend,
    radius       = 3,
    color        = {blue!40, orange!50, green!40, red!40, purple!30},
    explode      = {0, 0.1, 0, 0, 0},
  ]{
    35/[Kategori A],
    25/[Kategori B],
    20/[Kategori C],
    12/[Kategori D],
    8/[Lain-lain]
  }
\end{tikzpicture}
\caption{[Judul Diagram --- sumber: [PERLU VERIFIKASI]]}
\label{fig:komposisi}
\end{figure}

%------------------------------------------------------------
%  3. RADAR / SPIDER CHART — Perbandingan Skor Hukum
%     Perintah: /generate-viz radar-perbandingan
%------------------------------------------------------------
\begin{figure}[h]
\centering
\begin{tikzpicture}
\begin{polaraxis}[
  width            = 7cm,
  xtick            = {0, 60, 120, 180, 240, 300},
  xticklabels      = {Asas Kepastian, Asas Keadilan, Asas Kemanfaatan,
                      Efektivitas Norma, Aksesibilitas, Konsistensi},
  ymin=0, ymax=100,
  ytick            = {20, 40, 60, 80, 100},
  yticklabel style = {font=\tiny},
  legend style     = {at={(1.3,1)}, anchor=north west},
]
  % Ganti nilai-nilai (0-100) dengan skor riil Anda
  \addplot+[mark=*, thick, fill=blue!20, fill opacity=0.4]
    coordinates {(0,80)(60,65)(120,75)(180,55)(240,70)(300,60)(360,80)};
  \addplot+[mark=square*, thick, fill=orange!20, fill opacity=0.4]
    coordinates {(0,60)(60,75)(120,55)(180,80)(240,50)(300,70)(360,60)};
  \legend{[Sistem A], [Sistem B]}
\end{polaraxis}
\end{tikzpicture}
\caption{[Judul Radar --- sumber: [PERLU VERIFIKASI]]}
\label{fig:radar}
\end{figure}

%------------------------------------------------------------
%  4. TABEL PERBANDINGAN HUKUM (booktabs)
%     Perintah: /generate-viz tabel-perbandingan
%------------------------------------------------------------
\begin{table}[h]
\centering
\caption{Perbandingan [Topik] antar [Yurisdiksi/Regulasi]}
\label{tab:perbandingan}
\begin{tabular}{@{}p{3.5cm} p{3.5cm} p{3.5cm} p{3.5cm}@{}}
\toprule
\textbf{Aspek} & \textbf{[Sistem/Negara A]} & \textbf{[Sistem/Negara B]} & \textbf{[Sistem/Negara C]} \\
\midrule
Dasar Hukum      & [UU/Pasal] & [UU/Pasal] & [UU/Pasal] \\
\addlinespace
Asas yang Dianut & [Asas]     & [Asas]     & [Asas]     \\
\addlinespace
Teori Dominan    & [Teori]    & [Teori]    & [Teori]    \\
\addlinespace
Doktrin Rujukan  & [Doktrin]  & [Doktrin]  & [Doktrin]  \\
\addlinespace
Efektivitas      & [Uraian]   & [Uraian]   & [Uraian]   \\
\addlinespace
Catatan          & [Catatan]  & [Catatan]  & [Catatan]  \\
\bottomrule
\end{tabular}
\vspace{0.3em}
\footnotesize Sumber: [PERLU VERIFIKASI]
\end{table}

\end{document}
